\chapter{Trigonométrie dans le triangle rectangle}
\label{corriges:15}
\section*{Exercice 15.1 -- Les trois rapports}
Par rapport à l'angle en $B$, le côté opposé est $AC=6$, l'adjacent $AB=8$ et l'hypoténuse $BC=10$. Donc $\sin\alpha=0{,}6$, $\cos\alpha=0{,}8$ et $\tan\alpha=0{,}75$. En mode degrés, $\alpha=\arcsin(0{,}6)\approx36{,}87\degre$, soit $37\degre$.

\section*{Exercice 15.2 -- Hauteur et recul}
La hauteur est opposée à l'angle : $h=4\sin70\degre\approx3{,}76\unite m$. Le recul est adjacent : $d=4\cos70\degre\approx1{,}37\unite m$. On peut contrôler $h^2+d^2\approx16$, l'écart éventuel provenant des arrondis.

\section*{Exercice 15.3 -- Visée d'un sommet}
On suppose le terrain horizontal et le bâtiment vertical. La différence de hauteur entre les yeux et le sommet vaut $15\tan28\degre\approx7{,}9756\unite m$. Il faut ajouter la hauteur des yeux : $H\approx9{,}48\unite m$. Oublier les $1{,}50\unite m$ donnerait la hauteur relative aux yeux, pas au sol.

\section*{Exercice 15.4 -- Identité et pente}
Avec les côtés $o,a,c$, Pythagore donne $o^2+a^2=c^2$. Diviser par $c^2$ donne $(o/c)^2+(a/c)^2=1$, soit l'identité.
La pente signifie dénivelé/distance horizontale $=0{,}24$, donc $\tan\alpha=0{,}24$. L'angle vaut $\arctan(0{,}24)\approx13{,}5\degre$, pas $24\degre$.
